\section{The ramified prime and the selected local tangent component}
\label{sec:cubic-local-tangent}

The weaker rigidity question concerns the projection of actual tangent
cocycles, rather than every Bloch--Kato class of the mixed coefficient.
The following calculations determine what the fixed inertia condition
at $269$ and the crystalline condition at the coefficient prime impose.
They retain the separate question of membership in the selected local
deformation component.

\begin{proposition}[AC-G47: local mixed cohomology at $269$]
\label{AC-G47}
Let $F=K_{(269)}$, where $K=\mathbb Q(\omega)$, and let
$\ell\ne269$.  Let $L/\mathbb Q_\ell$ be finite and contain
$\mu_3$.  Suppose $V$ is an unramified two-dimensional
$L$-representation of $G_F$, and $\chi$ is the nontrivial
tame cubic character from AC-K04.  For
$B=\chi V$, one has
\[
             B^{I_F}=0,\qquad H^1(F,B)=0.             \tag{G47.1}
\]
The same holds for $B^\vee$, and therefore for the four-dimensional
off-diagonal coefficient of the cubic Type II point.
If $\ell\ne3,269$, the same assertions hold for any stable
integral lattice $T\subset B$ and its quotient $B/T$.

In particular, at a point whose actual inertia action is
$\chi I_2\oplus1$ on its endpoint block and middle line,
every global first-order mixed
cohomology class is automatically locally zero at $269$.
Its local first-order deformation is conjugate to that fixed
inertia action.  Here the hypothesis is on the inertia representation
itself, not only its Weil--Deligne semisimplification.
For $\ell\ne3,269$, it also satisfies the
literal condition that the deformation be the identity on
inertia intersected with the residual kernel.  Thus that
condition imposes no additional annihilation of the global
mixed cohomology class at first order.

At the exceptional coefficient prime $\ell=3$, retaining
that literal residual-kernel kill rule would instead exclude
the intended cubic seed itself: its reduction is trivial on
inertia at $269$, while its characteristic-zero restriction
is not trivial.
\end{proposition}
\begin{proof}
We give a direct cocycle proof, including the integral distinction.
Let $M$ be $B$, $B^\vee$, or one of the indicated lattice or
torsion modules.  Wild inertia $P_F$ is pro-$269$ and acts
trivially.  The restriction of a continuous cocycle to $P_F$
is consequently a continuous homomorphism to the additive
group of $M$.  Its image is compact.  In the coefficient-field
or lattice case a compact additive subgroup is pro-$\ell$;
in the discrete torsion case the image is a finite $\ell$-group.
It is also a quotient of a pro-$269$ group, hence is zero.
The cocycle therefore factors on inertia through its tame
quotient $I_F/P_F$, which is abelian.

Choose $\tau$ in this tame group with $\chi(\tau)\ne1$.
On $B$ its action is the scalar $\zeta=\chi(\tau)$;
on $B^\vee$ it is $\zeta^{-1}$.
The scalar $\zeta-1$ is invertible in $L$.
If $\ell\ne3$, it is a unit of $\mathcal O_L$, so is
also an automorphism of the lattice and torsion module.
For any cocycle $z$ and any tame inertia element $s$,
the equality $\tau s=s\tau$ gives
\[
             (\zeta-1)z(s)=(s-1)z(\tau).             \tag{G47.2}
\]
Thus $z|_{I_F}$ is the coboundary $s\mapsto X-sX$ with
$X=-(\zeta-1)^{-1}z(\tau)$, with the inverse character
used in the dual case.  This $X$ lies in the integral or
torsion module when $\ell\ne3$.  Moreover $M^{I_F}=0$
because $\tau-1$ is invertible.

Subtract that coboundary from a cocycle of $G_F$ and denote
the result by $z_0$.  It is zero on $I_F$.  For $g\in G_F$
and $s\in I_F$, the two expressions for $z_0(gs)$, using
$gs=(gsg^{-1})g$, show that $z_0(g)$ is fixed by every
element of $I_F$.  Therefore $z_0(g)=0$ for every $g$.
This proves (G47.1) and its stated integral variants without
an appeal to a dimension formula.
The same argument applies to their direct sum.  Since $K$ is
unramified at $269$, its local extension of $\mathbb Q_{269}$
has the same inertia group.  That direct sum therefore also
has zero $H^1$ over $\mathbb Q_{269}$: the proof after
(G47.2) uses only the vanishing of inertia cohomology and of
inertia invariants, regardless of the Frobenius action on the
two summands.

Now let $r$ be the cubic Type II point and let $z$ take values
in its off-diagonal adjoint coefficient.  The first-order lift
over $L[\varepsilon]/(\varepsilon^2)$ is
\[
                  r_z(g)=(I+\varepsilon z(g))r(g),    \tag{G47.3}
\]
with the second semidirect component unchanged.
The cocycle identity is exactly the homomorphism condition.
By (G47.1), locally at $269$ there is a mixed matrix $X$ with
$z(g)=X-\operatorname{Ad}(r(g))X$.  Consequently
\[
           r_z(g)=(I+\varepsilon X)r(g)
                              (I-\varepsilon X).     \tag{G47.4}
\]
This preserves the inertia conjugacy class and its characteristic
polynomials.  When $\ell\ne3,269$, reduction is injective on
$\mu_3$.  Hence the residual inertia kernel is exactly
$\ker(\chi|_{I_F})$.  On that kernel $r(g)=I$ and its
adjoint action is the identity, so $z(g)=0$ and $r_z(g)=I$.
The assertion is about each local restriction of a global class;
it does not conclude that the global class is zero.

At $\ell=3$, by contrast, both nontrivial cubic roots reduce
to $1$.  Thus every inertia element is in the residual kernel
of the cubic seed.  The rule $r(g)=I$ for that entire kernel
contradicts $r|_{I_F}=\chi\oplus1\oplus\chi$.
This is a failure to contain the proposed basepoint, not a
subtlety about tangent spaces.  Although (G47.1) for coefficient
fields still holds at $\ell=3$, multiplication by $\zeta-1$
is then not an integral automorphism, and the integral assertions
above were explicitly restricted to $\ell\ne3$.

The exact source rule is
\path{Development/AutomorphicArithmetic/AutomorphicVarieties/UnitaryThreeGroups/RhoKApply.lean},
\texttt{Dat.Kill}, lines 61--62, together with
\texttt{Dat.IsPt.kill}, line 82.
In \path{UnitaryThreeGroups/TanCocIsBK.lean},
\texttt{tanCoc\_isBK}, the away-from-$\ell$ part first uses
that rule on inertia intersected with $\ker\overline r$,
then the private finite-index averaging lemma
\texttt{tanBK\_avg} to obtain an inertia coboundary.
The calculation above identifies its sharper effect on the
particular cubic mixed coefficient.
\end{proof}

\begin{proposition}[AC-G47-crystalline: the exact infinitesimal local condition]
\label{AC-G47-crystalline}
Let $r$ be a continuous coefficient-field representation in
the source local semidirect group, and suppose it is crystalline
after restriction to the finite-index subgroup $H_L$ used in
the local datum.  Let $z$ be an adjoint cocycle and $r_z$ its
first-order deformation.  Then
\[
 r_z|_{H_L}\text{ is crystalline}
 \quad\Longleftrightarrow\quad
 [z]\in H^1_f(G_{\mathbb Q_\ell},\operatorname{ad}r).  \tag{G47.5}
\]
Here the finite local condition is the source's period-ring
condition and the finite-index descent convention is retained.
In particular this crystalline condition means that a class
is a coboundary after tensoring with $B_{\mathrm{cris}}$;
it does not assert that it is already a coboundary over the
coefficient field.

For the actual source ideal $I_t$, the verified implication is
\[
       \text{a first-order point on }I_t
                  \quad\Longrightarrow\quad [z]\in H^1_f.
                                                               \tag{G47.6}
\]
Replacing it by an equivalence at a non-base Type II point
requires an additional comparison with the selected component.
The supplied statements cited below establish the converse
formal comparison at the distinguished basepoint, not at
every such Type II point.  No assertion is made here that
the actual global tangent projection equals the entire mixed
Bloch--Kato group or that it is zero.
\end{proposition}
\begin{proof}
First restrict to $H_L$, on which the second semidirect
component is trivial and the matrix component of $r$ is a
linear representation.  Over the coefficient field its
first-order deformation is an extension of $r|_{H_L}$ by
itself.  Tensoring this extension with
$B_{\mathrm{cris}}$ makes it split exactly when the adjoint
cocycle $z$ becomes a coboundary in
$\operatorname{ad}r\otimes B_{\mathrm{cris}}$.
If the extension is crystalline, its crystalline periods have
twice the dimension of those of $r$; the resulting map onto
the periods of the quotient is surjective.  A lift of a basis
of those quotient periods gives an equivariant splitting after
tensoring with $B_{\mathrm{cris}}$.  Conversely such a splitting
exhibits the deformed representation as two copies of the
original one over the period ring, so it is crystalline.
This proves the equivalence after restriction to $H_L$.

The source subgroup $H_L=\ker\eta\cap\ker\tau$ is open
and normal.  If a period-ring-valued cocycle is a coboundary
on $H_L$, restriction--corestriction, or the average over its
finite cosets, shows that its global cohomology class is killed
by $[G_{\mathbb Q_\ell}:H_L]$.  This integer is invertible
in the coefficient field.  The class is therefore a global
period-ring coboundary.  Restriction gives the converse.
This is exactly the local finite condition in (G47.5).

The supplied proof verifies this argument for the actual
semidirect representation, including its inverse-transpose
action, in
\path{Development/AutomorphicArithmetic/AutomorphicVarieties/PadicHodgeTheory/DualCr.lean},
\texttt{dualCr}.  Its conclusion is literally
\texttt{IsLCr r' iff IsLCr r and gamma in finiteCocycles}.
The quantified base representation $r$ there is arbitrary;
it is not restricted to the diagonal basepoint.
The proof uses \texttt{deformCr} and finite-index descent
\texttt{z1fDesc}.

The separate component definition is in the same directory's
\path{IsInertiaInfiniteCyclotomic.lean}:
\texttt{Dat3.ordIdeal} fixes the stated inertia polynomial
and ordered-product equations;
\texttt{IY} saturates by \texttt{del};
\texttt{Icr} is its radical;
\texttt{primesAt0} selects minimal primes passing through
the distinguished basepoints; and
\texttt{It} is their intersection.
Thus component membership is additional structure in the
definition, even when all its points are crystalline.

In \path{PadicHodgeTheory/RtDual.lean},
\texttt{rtDual} applies \texttt{rtPoints} and \texttt{dualCr}
to an actual map killing $I_t$, proving (G47.6).
The potential converse condition \texttt{CrPts}, in
\texttt{IsInertiaInfiniteCyclotomic.lean}, is formulated with
\texttt{InDcr}; that abbreviation is
\texttt{InDcrAt} at \texttt{D.rz E}, the distinguished
diagonal representation.  It is not a universal statement
about all coefficient-field points of $I_t$.
Correspondingly \path{PadicHodgeTheory/RtCompl.lean},
\texttt{rtCompl}, assumes \texttt{IsBase}; it identifies
\texttt{primesAt0} with the singleton $\{I_t\}$ and constructs
a formal power-series chart in twelve variables there.
\texttt{RtChart.lean}, \texttt{rtChart}, retains the same
basepoint hypothesis.
These exact domains prove the asserted scope distinction.
They neither supply nor rule out a new selected-component
tangent comparison at the Type II points needed for the
modified cubic global deformation problem.
\end{proof}
